
\documentclass{article}
\usepackage{colortbl}
\usepackage{makecell}
\usepackage{multirow}
\usepackage{supertabular}

\begin{document}

\newcounter{utterance}

\centering \large Interaction Transcript for game `bbh', experiment `object\_counting', episode 129 with reflexion{-}llama{-}3.1{-}8b{-}instruct.
\vspace{24pt}

{ \footnotesize  \setcounter{utterance}{1}
\setlength{\tabcolsep}{0pt}
\begin{supertabular}{c@{$\;$}|p{.15\linewidth}@{}p{.15\linewidth}p{.15\linewidth}p{.15\linewidth}p{.15\linewidth}p{.15\linewidth}}
    \# & \multicolumn{2}{c}{Player} && \multicolumn{2}{c}{Game Master} \\
    \hline

    \theutterance \stepcounter{utterance}  
    & & & \multicolumn{4}{p{0.6\linewidth}}{
        \cellcolor[rgb]{0.9,0.9,0.9}{
            \makecell[{{p{\linewidth}}}]{
                \texttt{\tiny{[P1$\langle$GM]}}
                \texttt{You have attempted this same task before and did not succeed. These are the plans you wrote after those attempts.} \\
\texttt{Trial \#1: New plan: I failed to count the fruits in the last question because I didn't consider the possibility of a single fruit being mentioned multiple times. In the question "I have two grapes, and a plum. How many fruits do I have?", I should have counted the two grapes as two separate fruits, in addition to the plum. Next time, I will make sure to count each instance of a fruit individually, regardless of whether it's mentioned multiple times in the question.} \\
\\ 
\texttt{Trial \#2: New plan: In my previous attempts, I only considered the quantity of each type of fruit mentioned in the question, but I didn't account for the possibility of a single fruit being mentioned multiple times. However, in the question "I have two grapes, and a plum. How many fruits do I have?", I correctly counted the two grapes as two separate fruits, in addition to the plum. This suggests that my initial plan was correct, but I should have applied it consistently. Next time, I will make sure to count each instance of a fruit individually, regardless of whether it's mentioned multiple times in the question, and I will also consider the quantity of each type of fruit mentioned.} \\
\\ 
\texttt{Trial \#3: New plan: In my previous attempts, I correctly identified the need to count each instance of a fruit individually, regardless of whether it's mentioned multiple times in the question. However, I failed to apply this strategy consistently. In the question "I have two grapes, and a plum. How many fruits do I have?", I correctly counted the two grapes as two separate fruits, in addition to the plum. This suggests that my initial plan was correct, but I should have applied it consistently. Next time, I will make sure to count each instance of a fruit individually, regardless of whether it's mentioned multiple times in the question, and I will also consider the quantity of each type of fruit mentioned. To achieve this, I will create a mental or written list of each fruit mentioned in the question, and then count each instance of each fruit individually. This will ensure that I accurately count the total number of fruits.} \\
\\ 
\\ 
\texttt{Questions that involve enumerating objects and asking the model to count them.} \\
\\ 
\texttt{Q: I have a blackberry, a clarinet, a nectarine, a plum, a strawberry, a banana, a flute, an orange, and a violin. How many fruits do I have?} \\
\texttt{A: 6} \\
\\ 
\texttt{Q: I have an orange, a raspberry, two peaches, a blackberry, an apple, a grape, a nectarine, and three plums. How many fruits do I have?} \\
\texttt{A: 11} \\
\\ 
\texttt{Q: I have a lettuce head, a head of broccoli, an onion, a stalk of celery, two carrots, a garlic, and a yam. How many vegetables do I have?} \\
\texttt{A: 8} \\
\\ 
\texttt{Q: I have two grapes, and a plum. How many fruits do I have?} \\
\texttt{A:} \\
            }
        }
    }
     \\ \\

    \theutterance \stepcounter{utterance}  
    & \multicolumn{4}{p{0.6\linewidth}}{
        \cellcolor[rgb]{0.561,0.737,0.561}{
            \makecell[{{p{\linewidth}}}]{
                \texttt{\tiny{[P1$\rangle$GM]}}
                \texttt{Based on the plans you've written after previous attempts, I will apply the strategy of counting each instance of a fruit individually, regardless of whether it's mentioned multiple times in the question, and considering the quantity of each type of fruit mentioned.} \\
\\ 
\texttt{For the question "I have two grapes, and a plum. How many fruits do I have?", I will create a mental or written list of each fruit mentioned in the question:} \\
\\ 
\texttt{1. Grape (2 instances)} \\
\texttt{2. Plum (1 instance)} \\
\\ 
\texttt{Now, I will count each instance of each fruit individually:} \\
\\ 
\texttt{{-} Grape: 2} \\
\texttt{{-} Plum: 1} \\
\\ 
\texttt{Total number of fruits: 2 (grapes) + 1 (plum) = 3} \\
\\ 
\texttt{Therefore, the answer to the question is: 3} \\
            }
        }
    }
    & & \\ \\

    \theutterance \stepcounter{utterance}  
    & & & \multicolumn{2}{p{0.3\linewidth}}{
        \cellcolor[rgb]{0.9,0.9,0.9}{
            \makecell[{{p{\linewidth}}}]{
                \texttt{\tiny{[GM$|$GM]}}
                \texttt{Target: 3} \\
            }
        }
    }
    & & \\ \\

    \theutterance \stepcounter{utterance}  
    & & & \multicolumn{2}{p{0.3\linewidth}}{
        \cellcolor[rgb]{0.9,0.9,0.9}{
            \makecell[{{p{\linewidth}}}]{
                \texttt{\tiny{[GM$|$GM]}}
                \texttt{game\_result = LOSE} \\
            }
        }
    }
    & & \\ \\

    \theutterance \stepcounter{utterance}  
    & & \multicolumn{4}{p{0.6\linewidth}}{
        \cellcolor[rgb]{0.95,0.95,0.95}{
            \makecell[{{p{\linewidth}}}]{
                \texttt{\tiny{[P1$\leftrightarrow$P]}}
                \texttt{Post{-}episode Reflexion (not sent to GM):} \\
\texttt{New plan: In my previous attempts, I correctly identified the need to count each instance of a fruit individually, regardless of whether it's mentioned multiple times in the question. However, I failed to apply this strategy consistently. To achieve this, I will create a mental or written list of each fruit mentioned in the question, and then count each instance of each fruit individually. This will ensure that I accurately count the total number of fruits. To further improve my strategy, I will also consider the possibility of a single fruit being mentioned multiple times in the question, and I will make sure to count each instance of that fruit separately. For example, if the question is "I have two apples, and two bananas. How many fruits do I have?", I will create a list of each fruit mentioned in the question:} \\
\\ 
\texttt{1. Apple (2 instances)} \\
\texttt{2. Banana (2 instances)} \\
\\ 
\texttt{Then, I will count each instance of each fruit individually:} \\
\\ 
\texttt{{-} Apple: 2} \\
\texttt{{-} Banana: 2} \\
\\ 
\texttt{Total number of fruits: 2 (apples) + 2 (bananas) = 4} \\
\\ 
\texttt{By following this strategy, I will be able to accurately count the total number of fruits in each question, regardless of how they are presented.} \\
            }
        }
    }
    & \\ \\

\end{supertabular}
}

\end{document}
